\subsection{The generalized near-group ring $R^\prime(4)$}
\label{sec:frontier-155}
Write the basis of $\RR{4}{155}$ as $1,h,g,hg,X,Y$, where the
invertibles form $C_2^2$, $h$ fixes $X,Y$, and $g$ exchanges them.
The ring belongs to the family
\begin{equation}\label{eq:nonselfdual-neargroup}
 Y=X^*,\qquad XY=1+h+m(X+Y),\qquad X^2=g+hg+m(X+Y),
\end{equation}
with $m=4$. The following uniform argument eliminates signed-table
search for this family.

\begin{proposition}\label{prop:nonselfdual-neargroup-spherical}
Every complex fusion category with rules
\eqref{eq:nonselfdual-neargroup}, for $m\ge0$, is spherical and has a
pseudo-unitary Galois conjugate.
\end{proposition}
\begin{proof}
In the sphericalization choose the dimension-one lifts $k_+$ of the
invertibles, and choose $X_+$ and $Y_+=g_+X_+$. Their common nonzero
dimension $\delta$ forces
$h_+X_+=X_+$, $h_+Y_+=Y_+$ and $X_+^*=Y_+$.
The invertible part of the self-dual product $X_+Y_+$ is $1_++h_+$,
by reciprocity; hence
\[
 X_+Y_+=1_++h_++s(X_++Y_+)+t(X_-+Y_-),\qquad s+t=m.
\]
Thus $\delta^2=2+r\delta$ with $r=2(s-t)$; changing both lifts if
necessary gives $0\le r\le2m$. The squared-norm global dimension and
the Frobenius--Perron dimension are
\[
 d_r=8+r^2\pm r\sqrt{r^2+8},\qquad
 D_m=8+4m^2+2m\sqrt{4m^2+8}.
\]
The ratio $d_r/D_m$ is an algebraic integer
\cite[Proposition~8.22]{ENO}. For nonzero even $r$, $r^2+8$ is not a
square: $(v-r)(v+r)=8$ would force $r=1$. The same holds for $2m$.
When $m>0$, the compositum admits an automorphism negating both
quadratic roots (or only the denominator root when $r=0$).
Multiplying the ratio by its image gives the positive rational
algebraic integer
\[
 \frac{r^2+8}{4m^2+8}\le1.
\]
It equals one, so $r=2m$ and $t=0$. For $m=0$ this holds directly.
The lift $g_+$ commutes with $X_+,Y_+$, since both orders have the same
underlying simple and nonzero dimension. Therefore
$X_+^2=Y_+^2=g_+X_+Y_+$, and all products of the chosen lifts close
among them. They form a full spherical fusion subcategory projecting
tensor-equivalently onto the original category. Its noninvertible
dimension is $m\pm\sqrt{m^2+2}$; conjugating to the positive root
gives pseudo-unitarity.
\end{proof}

For $m=4$ the dimensions are
\[
 d=4+3\sqrt2,\qquad D=72+48\sqrt2.
\]
There is no nontrivial faithful grading: the adjoint subring is the
whole ring. The pointed subrings are its only proper nontrivial fusion
subrings. Thus neither a grading nor rank-three subring descent resolves
this case.

The retained central-induction certificate gives a finite necessary
domain of five profiles, with center ranks $46,48,52,54,60$.
For clarity, its categorical reductions are as follows. Pointed
subcategory projection for all three involutions leaves one low frame
and indicators $(\nu_2(X),\nu_2(Y),\nu_2(h),\nu_2(g),\nu_2(hg))
=(0,0,1,-1,-1)$. The $1+h$ module action and stabilizer-rotation
criterion constrain the third indicators. Exact Gram, Galois, word
and balancing equations then force central lifts of $h,g,hg$.
The lift of $h$ is fermionic: if it were bosonic, de-equivariantization
would split the fixed object $X$ into dimensions $d/2$, whose norm
is $-1/2$, contrary to integrality. A lift of $g$ generates a semion
subcategory of the center, which splits off from the center.

In its modular centralizer the restricted adjoint algebra is
$1+A_1+A_2$. For its restriction matrix $V$, adjunction gives
$V^{\mathsf T}V=I+N_{A_1}+N_{A_2}$. The two low simples have opposite
monodromy charges against the fermion. Consequently their multiplication
matrices are the same-charge and opposite-charge blocks of
$V^{\mathsf T}V-I$. Exhausting charge counts, duality, pointed actions,
restriction products and ribbon balancing leaves the six profiles,
each with a forced pair of low fusion matrices. For the rank-$44$ profile, the height-six full-character system of
Proposition~\ref{prop:round3-full-character}, in $\QQ(\zeta_{24})$,
has no solution with integral twist-type and rotation-eigenvalue counts.
The exact certificate gives $49$ integral-bound inferences followed
by one separating ray, using the bound argument in the proof of
Theorem~\ref{thm:1204-excluded}. Pointed restriction columns use their
proved twists $\pm1$ or $\pm i$; the other columns retain the complete
joint trace domain. All surviving character assignments differ only
by permutations of identical low columns, and no third-indicator value
is imposed. This removes the rank-$44$ profile.

The five remaining profiles are necessary conditions. The semion splitting of the center does
not assert a tensor-product splitting of the original category.

\begin{proposition}[No braided Morita representative]
\label{prop:155-no-braided-morita}
A categorification of $\RR{4}{155}$ is not Morita equivalent to any
braided fusion category.
\end{proposition}
\begin{proof}
Work with a pseudo-unitary Galois representative. Since the adjoint
subcategory is the whole category, the forgetful map on central
invertibles is injective. The central-induction certificate gives all four
central lifts, with twists $1,-1,i,-i$, up to complex conjugation.
Thus the center has no nonidentity invertible of twist one.

If a braided category $\mathcal M$ were Morita equivalent, it would
also be pseudo-unitary. The equivalence of centers preserves their
canonical positive ribbon structures. Every central invertible
forgetting to the unit of $\mathcal M$ has twist one, so this
forgetful map is injective. The image $a$ of the central semion has
its canonical lift supplied by the braiding of $\mathcal M$; uniqueness
forces that lift to be the semion itself. Hence $\langle a\rangle$
is nondegenerate, and braided decomposition gives
$\mathcal M\simeq\mathrm{Semion}\boxtimes\mathcal M_0$
\cite{DGNO}. Its center then contains the nonidentity twist-one
invertible $(a,a)$ from
$Z(\mathrm{Semion})\simeq\mathrm{Semion}\boxtimes\mathrm{Semion}^{\rm rev}$,
a contradiction. Galois conjugation preserves braidings and Morita
contexts, so the conclusion applies to the original category.
\end{proof}
